%!TEX root = TQFTbook.tex
%%%%%%%%%%%%%%%%%%%%%%%%%%%%%%
\chapter{Category of cobordisms and definition of TQFT}\label{c:cobordisms}

In this chapter, we give the key definition of this book, namely the notion of a
TQFT, following work of Atiyah (see \ocite{atiyah}). To do that, we first define
the notion of a cobordism.

Unless specified otherwise, ``manifold'' means a compact oriented smooth
manifold, possibly with boundary. For a manifold $M$, we denote by $\del M$ its
boundary, with the orientation determined by the orientation of $M$. If $\del
M=\varnothing$, we say that $M$ is \emph{closed}. Morphisms between manifolds
are smooth maps; unless stated otherwise, ``diffeomorphism'' means an
orientation-preserving invertible smooth map.

We denote by $\ov{M}$ the manifold $M$ considered with the
opposite orientation.

The main references for this chapter are \ocite{kock} and
\ocite{milnor-h-cobordism}.

%%%%%%%%%%%%%%%%%%%
\section{Category $\nCob$ and definition of TQFT}\label{s:TQFT-def}
\begin{definition}\label{d:cobordism}\index{cobordism}
Let $N_0, N_1$ be closed $(n-1)$-dimensional manifolds. A \emph{cobordism}
from $N_0$ to $N_1$ is an $n$-dimensional manifold $M$ together with
an isomorphism
$$
  \ph\colon \ov{N_0}\sqcup N_1 \isoto \del M.
$$
We will refer to $N_0$ as the incoming boundary (or simply the in-boundary),
and to $N_1$ as the outgoing boundary (out-boundary).

Following \ocite{kock}, we will write $N_0\cob{M} N_1$ to indicate that $M$ is a
cobordism from $N_0$ to $N_1$.
\end{definition}

The picture below shows a 2-dimensional cobordism between a circle and a union
of two circles; this cobordism is usually called \emph{pair of pants}\index{pair of pants}.
\begin{center}
\begin{tikzpicture}
  \pic[scale =1.5] {pants};
\end{tikzpicture}
\end{center}

\begin{convention}
In our pictures, cobordisms go ``top to bottom'':
the in-boundary is at the top, and the out-boundary is at the bottom.
\end{convention}

\begin{definition}\label{d:cobordism-isomorphism}\index{cobordism!isomorphism}
  An isomorphism of two cobordisms $M_1, M_2\in \nCob(N_0,N_1)$ is a
  diffeomorphism $f\colon M_1\to M_2$ which makes this diagram commutative:
  \begin{equation}\label{e:cob-iso}
  \begin{tikzcd}[column sep=5mm, row sep=2mm]
    & M_1\arrow[dd, "f"]\\
    \ov{N_0}\sqcup N_1
    \arrow[ur]
    \arrow[dr]\\
    & M_2
  \end{tikzcd}
  \end{equation}
\end{definition}
Given two cobordisms $N_0\cob{M_1}N_1$, $N_1\cob{M_2}N_2$, it is natural to
ask if they can be glued together in a single cobordism $M=M_2\circ M_1$
$$
  N_0\cob{M_2\circ M_1}N_2.
$$
Clearly, as a manifold, $M$ should be obtained by gluing $M_1$ and $M_2$ along
the common boundary $N_1$: $M= M_1\cup_{N_1} M_2$. However, there is a problem:
it is easy to define $M$ as a topological space, but the smooth structure is not
unique. There are several ways to deal with this problem.
\begin{itemize}
  \item One can switch to the category of PL manifolds and cobordisms,
    where the problem does not arise: the PL structure on $M_1\cup_{N_1} M_2$
    is uniquely determined by the PL structures on $M_1$, $M_2$.
  \item One can consider \emph{collared}\index{collared manifold}
    manifolds: instead of
    isomorphism $\ov{N_0}\sqcup N_1\to \del M$,
    we can require in the definition of cobordism an isomorphism of
    $(\ov{N_0}\sqcup N_1)\times [0, \epsilon)$ with a neighborhood
    of $\del M$ in $M$.
\end{itemize}
However, the simplest way is to note that while the smooth structure is not
unique, different smooth structures give isomorphic cobordisms.

\begin{theorem}\label{t:gluing-unique}
  The smooth structure on $M=M_1\cup_{N_1} M_2$ which agrees with
  the given smooth structures on $M_1$, $M_2$, is unique up to isomorphism:
  if $\al, \al'$ are two different smooth structures on $M$, then $(M,\al)$,
  $(M,\al')$ are isomorphic as cobordisms.
\end{theorem}
You can find a proof of this in \ocite{kock}*{Theorem 1.3.12}.

\begin{definition}\label{d:ncob}\index{cobordism category}
  The cobordism category $\nCob$ is the category whose objects are closed
  $(n-1)$-dimensional manifolds, and $\Hom(N_0, N_1)$ is the set of
  isomorphism classes of cobordisms $N_0\cob{M}N_1$. The composition of
  cobordisms $N_0\cob{M_1}N_1$, $N_1\cob{M_2}N_2$
  is given by
  $$
    M_2\circ M_1 = M_1\cup_{N_1} M_2
  $$
  (by the previous theorem it is independent of the choice of smooth structure).
  The identity morphism $\id_N$ is the cylinder $N\times I$.
\end{definition}
Similarly one defines the PL version of cobordism category.

This category has a natural structure of a symmetric monoidal
category with respect to disjoint union.

\begin{definition}\label{d:tqft}\index{TQFT}
  An $n$-dimensional TQFT is a symmetric
  monoidal functor $Z\colon \nCob\to \vect$.
\end{definition}
This definition was first given in \ocite{atiyah} in slightly different
language.

In particular, this definition implies that
\begin{equation}\label{e:tqft-gluing0}
  Z( M_1\cup_{N_1} M_2) = Z(M_2) Z(M_1)
\end{equation}
which is one of the forms of the \emph{gluing axiom}\index{gluing axiom}.
An equivalent formulation of the gluing axiom will be given below, see
\leref{l:tqft-gluing}.


One can also generalize this definition, allowing the TQFT to take values in
other categories.

\begin{definition}\label{d:tqft-c}
  Let $\calC$ be a symmetric monoidal category. Then an
  $n$-dimensional TQFT with values in $\calC$ is a symmetric
    monoidal functor $Z\colon \nCob\to \calC$.
\end{definition}
%%%%%%%%%%%%%%%%%%%%%%%%%%%%
\section{1D TQFT}\label{s:1d-tqft}

Now that we have a definition of a TQFT, let us try to construct examples. We
begin with the simplest possible case, that of a 1-dimensional TQFT.

\begin{theorem}\label{t:1d-TQFT}
  A 1-dimensional TQFT is the following collection of data:
  \begin{itemize}
    \item A vector space $V_+=Z(\bullet^+)$
    \item A vector space $V_-=Z(\bullet^-)$
    \item A linear map $\ev= Z(\cup)\colon V_+\otimes V_-\to \kk$
    \item A linear map $\coev = Z(\cap)\colon \kk \to V_-\otimes V_+$
  \end{itemize}
  satisfying relations below.
  \begin{equation}\label{e:1d-tqft-rigidity}
    \begin{aligned}
      &
      \begin{tikzpicture}[vcenter]
        %\draw[->] (0,0)--(0,-1) arc (-180:0:0.25) arc (180:0:0.25) -- +(0,-1);
        \pic[->]{rigidity};
        \node at (1.5, -1) {$=$};
        \draw [->] (2,0)--(2,-2);
      \end{tikzpicture}
      \qquad
      (V_+\to V_+ \otimes V_- \otimes V_+ \to V_+) = \id_{V_+}\\[20pt]
      &
      \begin{tikzpicture}[vcenter]
        %\draw[->] (0,0)--(0,-1) arc (-180:0:0.25) arc (180:0:0.25) -- +(0,-1);
        \pic[<-, x=-1cm]{rigidity};
        \node at (0.5, -1) {$=$};
        \draw [<-] (1,0)--(1,-2);
      \end{tikzpicture}
      \qquad
      (V_-\to V_- \otimes V_+ \otimes V_- \to V_-) = \id_{V_-}
    \end{aligned}
  \end{equation}
\end{theorem}
\begin{proof}
  Since $Z$ is symmetric monoidal, $Z(N_1\sqcup N_2)\simeq Z(N_1)\otimes Z(N_2)$,
  so $Z$ on objects is determined by its values on connected 0-manifolds. The
  only connected 0-manifolds are the positively and negatively oriented points
  $\bullet^+, \bullet^-$, hence $Z$ on objects is determined by
  $V_\pm := Z(\bullet^\pm)$.

  Similarly, $Z$ on morphisms is determined by its values on connected
  cobordisms. Up to isomorphism, every connected oriented $1$-cobordism with
  non-empty boundary is one of the following:
  \begin{itemize}
    \item the intervals $\bullet^+\to\bullet^+$ and $\bullet^-\to\bullet^-$,
      on which $Z$ is the identity of $V_+$ resp.\ $V_-$;
    \item the \emph{cup}
      $\bullet^+\sqcup\bullet^-\to\varnothing$, on which $Z$ is a linear
      map $\ev\colon V_+\otimes V_-\to\kk$;
    \item the \emph{cap}
      $\varnothing\to\bullet^-\sqcup\bullet^+$, on which $Z$ is a linear
      map $\coev\colon \kk\to V_-\otimes V_+$.
  \end{itemize}
  The only connected closed $1$-cobordism is the circle $S^1$; its value
  $Z(S^1)\in\kk$ will turn out to be determined by the data above.

  It remains to identify the relations these data must satisfy. The snake
  cobordism on the left-hand side of the first equation of
  \eqref{e:1d-tqft-rigidity} is diffeomorphic (rel boundary) to the straight
  interval on the right-hand side; applying $Z$ to this diffeomorphism gives
  the first snake identity. The second is obtained analogously by reversing
  orientations.

  Conversely, given data $(V_+, V_-, \ev, \coev)$ satisfying
  \eqref{e:1d-tqft-rigidity}, one defines $Z$ on objects and connected
  cobordisms by the formulas above and extends to general objects and
  morphisms by tensor product and composition. The relations
  \eqref{e:1d-tqft-rigidity} ensure that the result is independent of the
  decomposition; with this, $Z$ is well-defined as a symmetric monoidal
  functor $\oneCob\to\vect$.
\end{proof}

Comparing this with definition of dual pair given in \seref{s:rigid}, we see
that this data is  exactly the data of a dual pair  $V_+, V_-$. Therefore, by
\leref{l:rigidity_vect} $V_+$, $V_-$ must be finite-dimensional, and $V_-\simeq
(V_+)^*$, so we get the following result.

\begin{theorem}\label{t:1d-tqft2}
  A one-dimensional TQFT is completely determined by the vector space
  $V=Z(\bullet^+)$ which must be \emph{finite-dimensional}. Conversely, every
  finite-dimensional vector space $V$ uniquely defines 1-dimensional TQFT. For
  such a TQFT, $Z(S^1)=\dim V$.
\end{theorem}

This result easily generalizes to TQFTs with values in an arbitrary
symmetric monoidal category.
\begin{theorem}\label{t:1d-tqft}
  If $Z\colon \oneCob\to \calC$ is a $1$-dimensional TQFT with values in a
  symmetric monoidal category $\calC$, then $A=Z(\bullet^+)$ is a rigid object;
  its dual is $Z(\bullet^-)$. Conversely, any rigid object $A\in \calC$ defines
  a one-dimensional TQFT such that $Z(\bullet^+)=A$.
\end{theorem}

%\begin{tikzpicture}
%  \begin{scope}[tikzdefault]
%  \draw[->] (-0.5,0) arc (-180:0:0.5);
%  \draw[-] (0.5,0) arc (0:180:0.5);
%  \end{scope}
%\end{tikzpicture}


%%%%%%%%%%%%%%%%%%%%%
\section{First properties of TQFTs}\label{s:first-properties}



Let us now try to establish some properties of general $n$-dimensional TQFTs.

\begin{theorem}\label{t:tqft-rigidity}
  Let $Z\colon \nCob\to \calC$ be an $n$-dimensional TQFT with values in a
  symmetric monoidal category $\calC$. Then for any closed $(n-1)$-dimensional
  manifold $N$, $Z(N)\in \calC$ is rigid, and $Z(\ov{N})$ is a dual of $Z(N)$.
\end{theorem}
The proof is the same as in the one-dimensional case.

In particular, for TQFTs with values in $\vect$, for any closed
$(n-1)$-dimensional manifold $N$, the vector space $Z(N)$ is finite-dimensional.

\begin{exercise}\label{ec:tqft-duality}
  Let $Z$ be an $n$-dimensional TQFT with values in $\vect$.
  Prove that then
  $Z(N\times S^1)=\dim Z(N)$ for any closed $(n-1)$-dimensional manifold $N$.
\end{exercise}

More generally, we have the following lemma.
\begin{lemma}\label{l:tqft-gluing}
Let $Z$ be an $n$-dimensional TQFT, and let $M$ be a $n$-dimensional
manifold with boundary. Let $N\subset M$ be a a smooth codimension one
submanifold and denote by $M_{\text{cut}}$ the manifold with boundary
obtained by cutting $M$ along $N$; thus,
$\del(M_{\text{cut}})\cong N\sqcup \ov{N}\sqcup \del M$,
see \firef{f:cutting}.

Prove that then
$$
  Z(M)=\tr_{Z(N)} (Z(M_{\text{cut}}))\in Z(\del M)
$$
where $\tr_{Z(N)}$ is the canonical pairing
$$
  Z(N)\otimes Z(\ov{N})\cong Z(N)\otimes Z(N)^*\to \kk
$$
\end{lemma}

\begin{figure}[ht]
  \centering

  \begin{tikzpicture}[scale=0.5]
    \begin{scope}[xshift=0cm]
      % first diagram code
      \node (31) at (-3.75, -1.5) {};
      \node (32) at (-3.75, 0.5) {};
      \node (33) at (3, 1.25) {};
      \node (34) at (-1.75, 0.5) {};
      \node (35) at (1.5, 0.75) {};
      \node (36) at (-1.25, 0.27) {};
      \node (37) at (1, 0.43) {};
      \node (38) at (5.5, 0.25) {};
      \node (41) at (5.5, -2.25) {};
      \node (42) at (3, -2.75) {};
      \node (43) at (-1.25, -2.75) {};
      \node (44) at (0, 2.62) {};
      \node (45) at (0, 1.07) {};
      \draw [bend left=60] (32.center) to (33.center);
      \draw [bend left, looseness=0.75] (31.center) to (32.center);
      \draw [bend right, looseness=1.25] (34.center) to (35.center);
      \draw [bend left=60, looseness=1.25] (36.center) to (37.center);
      \draw [bend right, looseness=0.75] (33.center) to (38.center);
      \draw [bend left=15, looseness=1.25] (38.center) to (41.center);
      \draw [bend right] (31.center) to (43.center);
      \draw [bend left=15, looseness=1.25] (43.center) to (42.center);
      \draw [bend right, looseness=0.75] (42.center) to (41.center);
      \draw [bend right=15, looseness=1.25] (38.center) to (41.center);
      \draw [dashed, bend right, looseness=0.50] (44.center) to (45.center);
      \draw [dashed, bend left=15] (44.center) to (45.center);
      \node[left] at ($(45.center)+(-0.20,0.60)$) {$N$};
      \node[left] at ($(45.center)+(+5.50,-1.50)$) {$\del M$};
      \node[below] at ($(45.center)+(0.50,-4.00)$) {$M$};

    \end{scope}

    \begin{scope}[xshift=11cm]
      % second diagram code
      \node (31) at (-3.5, -1.5) {};
      \node (32) at (-3.5, 0.5) {};
      \node (33) at (3.25, 1.25) {};
      \node (34) at (-1.5, 0.5) {};
      \node (35) at (1.75, 0.75) {};
      \node (36) at (-1, 0.27) {};
      \node (37) at (1.25, 0.43) {};
      \node (38) at (5.75, 0.25) {};
      \node (41) at (5.75, -2.25) {};
      \node (42) at (3.25, -2.75) {};
      \node (43) at (-1, -2.75) {};
      \node (44) at (0, 2.75) {};
      \node (45) at (0, 1) {};
      \node (46) at (0.5, 2.75) {};
      \node (47) at (0.5, 1) {};
      \draw [bend left, looseness=0.75] (31.center) to (32.center);
      \draw [bend right, looseness=1.25] (34.center) to (35.center);
      \draw [bend right, looseness=0.75] (33.center) to (38.center);
      \draw [bend left=15, looseness=1.25] (38.center) to (41.center);
      \draw [bend right] (31.center) to (43.center);
      \draw [bend left=15, looseness=1.25] (43.center) to (42.center);
      \draw [bend right, looseness=0.75] (42.center) to (41.center);
      \draw [bend right=15, looseness=1.25] (38.center) to (41.center);
      \draw [bend right, looseness=0.75] (44.center) to (45.center);
      \draw [bend left=15] (44.center) to (45.center);
      \draw [bend right] (33.center) to (46.center);
      \draw [bend left=45, looseness=0.75] (32.center) to (44.center);
      \draw [bend left, looseness=1.25] (36.center) to (45.center);
      \draw [bend right, looseness=1.25] (37.center) to (47.center);
      \draw [bend left=15, looseness=1.25] (46.center) to (47.center);
      \draw [bend right=15] (46.center) to (47.center);
      \node[left] at ($(45.center)+(-0.10,0.60)$) {$N$};
      \node[right] at ($(45.center)+(+0.50,0.60)$) {$N$};
      \node[left] at ($(45.center)+(+5.50,-1.50)$) {$\del M$};
      \node[below] at ($(45.center)+(0.50,-4.00)$) {$M_{cut}$};

    \end{scope}

  \end{tikzpicture}

  \caption{Cutting of a manifold $M$ along a codimension one submanifold $N$}
  \label{f:cutting}
\end{figure}


This lemma is another form of the gluing axiom for TQFT (compare with
\eqref{e:tqft-gluing0}.)

Note that unlike the one-dimensional case, in general we cannot claim that any
rigid object of $\calC$ defines a TQFT. The proper generalization of that
statement is the cobordism hypothesis, which we will discuss later.



Another property of TQFTs is functoriality. As mentioned before, it is
natural to expect that a diffeomorphism $\ph \colon N_0\to N_1$ between two
$(n-1)$-dimensional closed manifolds gives rise to an isomorphism
$Z(\ph)\colon Z(N_0)\to Z(N_1)$; however, we didn't include that as part of
the definition. It turns out it is not necessary.

Let $\ph \colon N_0\to N_1$ be an orientation-preserving diffeomorphism.
Define the cobordism
\begin{equation}\label{e:Cph}
  N_0\cob{C_\ph} N_1
\end{equation}
to be the cylinder $N_1\times I$, with the identification $\ov{N_0}\sqcup N_1
\to \del C_\ph$ given by $\ph\sqcup \id$.

\begin{lemma}\label{l:tqft-isotopy}
  Let $C_\ph$ be defined above. Then:
  \begin{enumerate}
    \item The isomorphism class of $C_\ph$ \textup{(}as a cobordism\textup{)}
    only depends on the
    isotopy class of $\ph$: if $\ph, \ph'$ are isotopic to each other, then
      $C_\ph\simeq C_{\ph'}$.
    \item Given diffeomorphisms $\ph_1\colon N_0\to N_1$, $\ph_2\colon N_1\to
      N_2$, we have an isomorphism of cobordisms
      $$
        C_{\ph_2}\circ C_{\ph_1}\simeq C_{\ph_2 \ph_1}.
      $$
  \end{enumerate}
\end{lemma}
Recall that $\ph, \ph'$ are called isotopic if there exists a one-parameter
family $\ph_t$ of smooth maps $N_0\to N_1$ such that $\ph_t$ is a diffeomorphism
for all $t$, and $\ph_0=\ph$, $\ph_1=\ph'$.

We leave the proof of this lemma as an easy exercise to the reader.

\begin{remark}
  In fact, there is a stronger statement: $C_\ph\simeq C_{\ph'}$ if and only if
  $\ph, \ph'$ are \emph{pseudo-isotopic}\index{pseudo-isotopic}, see
  \ocite{milnor-h-cobordism}*{Theorem 1.9}.
\end{remark}




As an immediate corollary, we get the following theorem.
\begin{theorem}\label{t:tqft-functoriality}
  Let $Z\colon \nCob\to \calC$ be an $n$-dimensional TQFT.
  \begin{enumerate}
    \item For any orientation-preserving diffeomorphism $\ph\colon N_0\to N_1$, define
      $Z(\ph)\colon Z(N_0)\to Z(N_1)$ by $Z(\ph)=Z(C_\ph)$. Then
      $Z(\ph_2\ph_1)= Z(\ph_2)\circ Z(\ph_1)$.
    \item If $\ph,\ph'$ are isotopic, then $Z(\ph)=Z(\ph')$.
  \end{enumerate}
\end{theorem}
\clearpage
